\chapter{Refraction of Light}\label{ch:g10:refraction}

Dip a straw into a glass of water: it looks snapped at the waterline.
Swimming pools are always deeper than they look, diamonds flash colored
fire from a colorless stone, and the sentence you are reading probably
crossed an ocean inside a glass thread thinner than a hair. One law
explains all four --- an equality between two sines, found by Snell in
1621 after fifteen centuries of physicists staring at tables of angles
without seeing the pattern.

\section{Reflection: a reminder}

In a homogeneous transparent medium, light travels in straight lines ---
the rays of \cref{ch:g10:light-spectra}. At the boundary between two
media, part of the light bounces back (\emph{reflection}) and part
crosses over (\emph{refraction}). Both obey exact geometric laws, stated
with one convention: angles are measured from the \emph{normal}, never
from the surface.

\begin{definition}[Normal and angle of incidence]\label{def:g10:refraction:angles}
At the point where a ray meets a surface, the
\emph{normal}\index{normal} is the line perpendicular to the surface at
that point. The \emph{angle of incidence}\index{angle of incidence} $i_1$
is the angle between the incoming ray and the normal; the angles of
reflection and of refraction are measured from the normal in the same
way.
\end{definition}

\begin{proposition}[Law of reflection]\label{prop:g10:refraction:reflection}
The reflected ray lies in the plane containing the incident ray and the
normal, on the other side of the normal, and the angle of reflection
equals the angle of incidence:
\[
i_1' = i_1 .
\]
\end{proposition}
\admitted

\begin{example}[Reading angles correctly]\label{ex:g10:refraction:mirror}
A laser beam strikes a mirror ``at $35^\circ$'' --- measured from the
mirror's surface, as beams are often described. The angle of incidence is
$90^\circ - 35^\circ = 55^\circ$, so the reflected ray also leaves at
$55^\circ$ from the normal, and the beam turns by
$180^\circ - 2 \times 55^\circ = 70^\circ$. Forgetting the
normal convention is the most common error of this chapter; it costs
$90^\circ$ every time.
\end{example}

\section{Refraction and the refractive index}

Hold a pencil in a half-filled glass: it looks broken at the surface.
Light from the underwater part changes direction as it leaves the water,
so the eye --- which assumes straight rays --- reconstructs the tip
somewhere it is not.

\begin{definition}[Refraction]\label{def:g10:refraction:refraction}
\emph{Refraction}\index{refraction} is the change of direction of light
as it crosses the boundary between two transparent media. The ray in the
second medium is the \emph{refracted ray}\index{refracted ray}, and its
angle from the normal is the angle of refraction $i_2$.
\end{definition}

Refraction happens because light has different speeds in different
media; each material is characterized by how much it slows light down.

\begin{definition}[Refractive index]\label{def:g10:refraction:index}
The \emph{refractive index}\index{refractive index} of a transparent
medium is the ratio
\[
n = \frac{c}{v},
\]
where $c = \qty{3.00e8}{m/s}$ is the speed of light in vacuum and $v$ its
speed in the medium. Since $v \leq c$, the index satisfies $n \geq 1$; it
has no unit. Some useful values:
\begin{center}
\begin{tabular}{lccccc}
medium & air & water & plexiglas & glass & diamond \\
$n$ & $1.0003$ & $1.33$ & $1.49$ & $1.50$ & $2.42$ \\
\end{tabular}
\end{center}
\end{definition}

\begin{example}[Speed of light in water]\label{ex:g10:refraction:speed}
In water, $v = \dfrac{c}{n} = \dfrac{\qty{3.00e8}{m/s}}{1.33}
= \qty{2.26e8}{m/s}$. In diamond, light crawls at
$\qty{3.00e8}{m/s}/2.42 = \qty{1.24e8}{m/s}$ --- less than half its
vacuum speed. For air, $n = 1.0003$ changes the speed by $0.03\%$: in
this chapter we treat air as vacuum and take $n_{\text{air}} = 1.00$.
\end{example}

\begin{theorem}[Snell's law of refraction]\label{thm:g10:refraction:snell}
\index{Snell's law}
The refracted ray lies in the plane of the incident ray and the normal,
on the other side of the normal, and the angles of incidence and
refraction satisfy
\[
n_1 \sin i_1 = n_2 \sin i_2 ,
\]
where $n_1$ and $n_2$ are the refractive indices of the two media. The
path is reversible: light traveling backwards along the refracted ray
exits along the incident ray.
\end{theorem}
\admitted

\begin{remark}\label{rem:g10:refraction:honest}
The laws of reflection and refraction are stated here as experimental
facts. They are honestly derived in the Year 1 volume, where light is
treated as a wave: both laws --- and the formula $n = c/v$ itself ---
then follow from the wave slowing down in the denser medium.
\end{remark}

\begin{omfigure}
\begin{tikzpicture}[scale=1.0, font=\small]
  \fill[blue!8] (-3.2,-2.1) rectangle (3.2,0);
  \draw[thick] (-3.2,0) -- (3.2,0);
  \draw[dashed] (0,-2.1) -- (0,2.3);
  \coordinate (O) at (0,0);
  \coordinate (I) at (-1.8,2.0);
  \coordinate (R) at (1.8,2.0);
  \coordinate (T) at (1.0,-1.73);
  \coordinate (N) at (0,2.3);
  \coordinate (Nd) at (0,-2.1);
  \draw[-Stealth, omDef, thick] (I) -- (O);
  \draw[-Stealth, omExo, thick] (O) -- (R);
  \draw[-Stealth, omThm, thick] (O) -- (T);
  \pic[draw, angle radius=9mm, "$i_1$", angle eccentricity=1.35] {angle = N--O--I};
  \pic[draw, angle radius=7mm, "$i_1'$", angle eccentricity=1.4] {angle = R--O--N};
  \pic[draw, angle radius=9mm, "$i_2$", angle eccentricity=1.35] {angle = Nd--O--T};
  \node[anchor=west] at (-3.1,0.45) {medium 1 ($n_1$)};
  \node[anchor=west] at (-3.1,-0.55) {medium 2 ($n_2$)};
\end{tikzpicture}

{\small Incident (blue), reflected (gray) and refracted (red) rays. All
angles are measured from the normal (dashed); entering the denser
medium ($n_2 > n_1$), the ray bends \emph{toward} the normal.}
\end{omfigure}

\begin{example}[Air to water]\label{ex:g10:refraction:airwater}
A ray hits a calm pond at $i_1 = 30^\circ$. With $n_1 = 1.00$ and
$n_2 = 1.33$:
\[
\sin i_2 = \frac{n_1 \sin i_1}{n_2} = \frac{\sin 30^\circ}{1.33}
= \frac{0.500}{1.33} = 0.376,
\qquad
i_2 = 22.1^\circ .
\]
The ray bends toward the normal, as it always does when entering a
slower medium ($n_2 > n_1$). Leaving the water, the same computation runs
backwards and the ray bends \emph{away} from the normal.
\end{example}

\begin{method}[Solving a refraction problem]\label{met:g10:refraction:method}
\begin{enumerate}
  \item Draw the surface, the normal, and the ray; identify $n_1$
        (medium of the incident ray) and $n_2$.
  \item Check that the given angles are measured from the normal;
        convert if they are measured from the surface.
  \item Write $n_1 \sin i_1 = n_2 \sin i_2$ and isolate the unknown
        sine; get the angle with the calculator's inverse sine, as in
        the mathematics volume.
  \item Sanity check: the ray bends toward the normal if $n_2 > n_1$,
        away from it if $n_2 < n_1$, and $\sin i_2$ must come out
        between $0$ and $1$.
\end{enumerate}
\end{method}

The experiment behind the law fits on a desk: a half-disc of plexiglas,
a protractor, a laser. Raw angle pairs $(i_1, i_2)$ look irregular ---
Ptolemy tabulated them around the year 150 and wrongly concluded that
$i_2$ was proportional to $i_1$. Plot sine against sine instead, and the
data straighten into a line through the origin of slope $n_1/n_2$: that
line \emph{is} Snell's law.

\begin{omfigure}
\begin{tikzpicture}
\begin{axis}[omaxis, width=7.2cm, height=5.4cm,
  xlabel={$\sin i_1$}, ylabel={$\sin i_2$},
  xmin=0, xmax=1.08, ymin=0, ymax=0.88,
  xtick={0,0.2,0.4,0.6,0.8,1.0}, ytick={0.2,0.4,0.6,0.8}]
  \addplot[omThm, thick, domain=0:1.02] {x/1.33};
  \addplot[only marks, mark=*, mark size=1.6pt, omDef] coordinates
    {(0.174,0.131) (0.342,0.257) (0.500,0.376) (0.643,0.483)
     (0.766,0.576) (0.866,0.651) (0.940,0.707)};
\end{axis}
\end{tikzpicture}

{\small Air-to-water measurements for $i_1 = 10^\circ$ to $70^\circ$:
plotted sine against sine, the data align on a line through the origin
of slope $1/1.33$.}
\end{omfigure}

\begin{proposition}[Apparent depth]\label{prop:g10:refraction:depth}
\index{apparent depth}
An object at depth $d$ under water, viewed from above at small angles,
appears to be at the \emph{apparent depth}
\[
d' = \frac{d}{n} ,
\]
where $n$ is the index of water. A pool of depth \qty{1.20}{m} looks
only $1.20/1.33 = \qty{0.90}{m}$ deep --- a classic cause of imprudent
dives.
\end{proposition}
\admitted

\begin{remark}\label{rem:g10:refraction:depthhonest}
The formula follows from Snell's law applied to the almost-vertical rays
reaching the eye; the derivation is done in the Year 1 volume. The
direction of the effect needs no formula: rays leaving the water bend
away from the normal, so the eye traces them back to a point higher
than the object.
\end{remark}

\section{Dispersion: one law per color}

A prism turns white sunlight into a rainbow, as the spectra of
\cref{ch:g10:light-spectra} showed. Refraction explains \emph{why}: the
refractive index of a medium is not quite the same for every color.

\begin{definition}[Dispersion]\label{def:g10:refraction:dispersion}
A medium is \emph{dispersive}\index{dispersion} when its refractive
index depends on the wavelength of the light. In glass and water, $n$ is
slightly larger for blue light (short wavelength) than for red light
(long wavelength) --- so blue bends more.
\end{definition}

\begin{example}[Two colors, one surface]\label{ex:g10:refraction:twocolors}
A beam of white light hits crown glass at $i_1 = 60^\circ$. For this
glass $n_{\text{red}} = 1.51$ and $n_{\text{blue}} = 1.53$. Snell's law,
once per color:
\[
\sin i_{\text{red}} = \frac{\sin 60^\circ}{1.51} = 0.574,
\quad i_{\text{red}} = 35.0^\circ ;
\qquad
\sin i_{\text{blue}} = \frac{\sin 60^\circ}{1.53} = 0.566,
\quad i_{\text{blue}} = 34.5^\circ .
\]
A single surface splits white light by half a degree --- invisible on a
window pane, but a prism refracts twice in the same rotational sense,
widens the fan, and a few meters away the colors are centimeters apart.
\end{example}

\begin{remark}\label{rem:g10:refraction:rainbow}
The rainbow is the same computation run inside a raindrop: sunlight
refracts entering the drop, reflects on its back wall, and refracts
again on the way out. Water's index runs from $1.331$ (red) to $1.343$
(blue), so each color exits at its own angle --- and no two observers
see the same rainbow.
\end{remark}

\section{Total internal reflection}

Snell's law contains a trap. Going from water toward air, the refracted
ray bends away from the normal: $\sin i_2 = 1.33 \sin i_1 > \sin i_1$.
Increase $i_1$, and the formula soon demands $\sin i_2 > 1$, which no
angle can deliver. Refraction then simply stops happening.

\begin{definition}[Critical angle and total internal reflection]\label{def:g10:refraction:critical}
Let light travel in a medium of index $n_1$ toward a medium of smaller
index $n_2 < n_1$. The \emph{critical angle}\index{critical angle}
$i_c$ is the angle of incidence for which the refracted ray grazes the
surface ($i_2 = 90^\circ$). For $i_1 > i_c$ no refracted ray exists: the
surface reflects all the light back into the first medium, following the
law of reflection. This is \emph{total internal
reflection}\index{total internal reflection}.
\end{definition}

\begin{proposition}[Critical angle formula]\label{prop:g10:refraction:criticalformula}
For $n_1 > n_2$,
\[
\sin i_c = \frac{n_2}{n_1} .
\]
\end{proposition}

\begin{proof}
Set $i_2 = 90^\circ$ in Snell's law: $n_1 \sin i_c = n_2 \sin 90^\circ
= n_2$. For $i_1 > i_c$, Snell's law would require
$\sin i_2 = \frac{n_1}{n_2}\sin i_1 > 1$: no refracted direction exists.
(That the light is then \emph{entirely} reflected --- not absorbed --- is
an experimental fact, confirmed by the wave theory of the Year 1
volume.)
\end{proof}

\begin{omfigure}
\begin{tikzpicture}[scale=0.92, font=\small]
  \fill[blue!8] (-3.6,-1.9) rectangle (3.6,0);
  \draw[thick] (-3.6,0) -- (3.6,0);
  \draw[dashed] (-2.2,-1.6) -- (-2.2,1.6);
  \draw[dashed] (0.3,-1.6) -- (0.3,1.6);
  \draw[dashed] (2.2,-1.6) -- (2.2,1.6);
  \draw[-Stealth, omDef, thick] (-2.95,-1.30) -- (-2.2,0);
  \draw[-Stealth, omThm, thick] (-2.2,0) -- (-1.20,1.12);
  \draw[-Stealth, omDef, thick] (-0.83,-0.99) -- (0.3,0);
  \draw[-Stealth, omThm, thick] (0.3,0) -- (1.75,0.06);
  \draw[-Stealth, omDef, thick] (0.90,-0.75) -- (2.2,0);
  \draw[-Stealth, omDef, thick] (2.2,0) -- (3.50,-0.75);
  \node at (-2.2,-1.85) {$i_1 < i_c$};
  \node at (0.3,-1.85) {$i_1 = i_c$};
  \node at (2.2,-1.85) {$i_1 > i_c$};
\end{tikzpicture}

{\small From water toward air, at increasing incidence: refraction away
from the normal, then a refracted ray grazing the surface at the
critical angle, then total internal reflection --- the surface has
become a perfect mirror.}
\end{omfigure}

\begin{example}[Three critical angles]\label{ex:g10:refraction:criticalvalues}
Toward air: for water, $\sin i_c = 1/1.33 = 0.752$, so
$i_c = 48.8^\circ$; for glass, $\sin i_c = 1/1.50$, so
$i_c = 41.8^\circ$; for diamond, $\sin i_c = 1/2.42 = 0.413$, so
$i_c = 24.4^\circ$. The higher the index, the narrower the escape cone:
in diamond, any ray more than $24.4^\circ$ off a facet's normal is
trapped --- the secret of its sparkle, dissected in
\cref{pb:g10:refraction:1}.
\end{example}

\section{Optical fibers}

Total internal reflection turns a glass thread into a pipe for light:
a ray entering nearly parallel to the axis strikes the wall far beyond
the critical angle, reflects without loss, and does so millions of times
per kilometer without leaking.

\begin{definition}[Optical fiber]\label{def:g10:refraction:fiber}
An \emph{optical fiber}\index{optical fiber} is a thin glass thread made
of a \emph{core}\index{core} of index $n_1$ wrapped in a
\emph{cladding}\index{cladding} of slightly smaller index $n_2 < n_1$.
Light injected at one end travels along the core by repeated total
internal reflection on the core--cladding boundary.
\end{definition}

\begin{omfigure}
\begin{tikzpicture}[scale=1.0, font=\small]
  \draw[omExo, thick] (0,0.85) -- (7,0.85);
  \draw[omExo, thick] (0,-0.85) -- (7,-0.85);
  \fill[blue!8] (0,-0.5) rectangle (7,0.5);
  \draw[thick] (0,0.5) -- (7,0.5);
  \draw[thick] (0,-0.5) -- (7,-0.5);
  \draw[-Stealth, omThm, thick] (0,-0.36) -- (1.0,0.5) -- (2.16,-0.5)
    -- (3.32,0.5) -- (4.48,-0.5) -- (5.64,0.5) -- (6.6,-0.33);
  \node[anchor=south] at (3.5,0.9) {cladding ($n_2 = 1.46$)};
  \node[anchor=north] at (3.5,-0.9) {core ($n_1 = 1.48$)};
\end{tikzpicture}

{\small Light guided along the core of a fiber by total internal
reflection (bounce angles greatly exaggerated: real rays graze the wall
at more than $80^\circ$ from the normal).}
\end{omfigure}

\begin{example}[A fiber by the numbers]\label{ex:g10:refraction:fibernumbers}
With $n_1 = 1.48$ and $n_2 = 1.46$,
\[
\sin i_c = \frac{1.46}{1.48} = 0.986,
\qquad i_c = 80.6^\circ :
\]
only rays within about $9^\circ$ of the axis are guided --- which is
fine, since that is exactly how laser light is injected. Inside the
core the light travels at $v = c/1.48 = \qty{2.03e8}{m/s}$: a pulse
crosses a \qty{100}{km} fiber in about half a millisecond.
\end{example}

\begin{remark}[Order of magnitude of data links]\label{rem:g10:refraction:datarates}
Submarine cables --- bundles of a few dozen fibers on the ocean floor
--- carry almost all intercontinental traffic. A single modern fiber
transports on the order of $10^{13}$ bits per second by sending many
wavelengths at once; a whole cable reaches $10^{14}$ bits per second,
with an amplifier every \qty{80}{km} or so. Satellites, for comparison,
handle a small fraction of a percent of the traffic.
\end{remark}

\section{Exercises}

\begin{exercise}[$\star$]\label{exo:g10:refraction:1}
A laser beam strikes a flat mirror at $25^\circ$ \emph{from the mirror's
surface}. Give the angle of incidence, the angle of reflection, and the
angle between the incident and reflected beams.
\end{exercise}

\begin{exercise}[$\star$]\label{exo:g10:refraction:2}
\begin{enumerate}
  \item Compute the speed of light in water ($n = 1.33$) and in diamond
        ($n = 2.42$).
  \item In a certain solid, light travels at \qty{1.97e8}{m/s}. Compute
        its refractive index and identify the material.
\end{enumerate}
\end{exercise}

\begin{exercise}[$\star$]\label{exo:g10:refraction:3}
A ray in air strikes a glass block ($n = 1.50$) at $i_1 = 40^\circ$.
Compute the angle of refraction. At what angle would a ray traveling
inside the glass have to hit the surface to exit into the air at
$40^\circ$?
\end{exercise}

\begin{exercise}[$\star$]\label{exo:g10:refraction:4}
Sunlight reaches the surface of a lake at $55^\circ$ from the vertical.
Compute the angle of the refracted ray, and state whether the ray bends
toward or away from the normal --- with the one-word reason.
\end{exercise}

\begin{exercise}[$\star$]\label{exo:g10:refraction:5}
To identify a clear liquid, a ray is sent onto its surface at
$i_1 = 45^\circ$; the refracted ray is measured at $i_2 = 32^\circ$.
Compute the refractive index and identify the liquid from the table of
\cref{def:g10:refraction:index}.
\end{exercise}

\begin{exercise}[$\star\star$]\label{exo:g10:refraction:6}
A half-disc of a transparent plastic gives the following measurements:
\begin{center}
\begin{tabular}{lccccc}
$i_1$ & $20^\circ$ & $30^\circ$ & $40^\circ$ & $50^\circ$ & $60^\circ$ \\
$i_2$ & $13.5^\circ$ & $19.5^\circ$ & $25.5^\circ$ & $31.0^\circ$
      & $35.5^\circ$ \\
\end{tabular}
\end{center}
\begin{enumerate}
  \item Compute $\dfrac{\sin i_1}{\sin i_2}$ for each pair. What do you
        observe?
  \item Deduce the refractive index of the plastic and identify it.
  \item Why is plotting $\sin i_2$ against $\sin i_1$ a better test of
        Snell's law than plotting $i_2$ against $i_1$?
\end{enumerate}
\end{exercise}

\begin{exercise}[$\star\star$]\label{exo:g10:refraction:7}
Compute the critical angle toward air for water, glass ($n = 1.50$) and
diamond. Then explain, with Snell's law, why there is no critical angle
for light going \emph{from} air \emph{into} water.
\end{exercise}

\begin{exercise}[$\star\star$]\label{exo:g10:refraction:8}
\begin{enumerate}
  \item A pool is \qty{2.0}{m} deep. What depth does a swimmer standing
        at the edge perceive?
  \item A heron eyes a fish swimming \qty{40}{cm} below the surface. At
        what depth does the fish appear, and should the heron strike
        above, at, or below the image it sees?
\end{enumerate}
\end{exercise}

\begin{exercise}[$\star\star$]\label{exo:g10:refraction:9}
A ray hits a window pane ($n = 1.50$, parallel faces) at $50^\circ$.
\begin{enumerate}
  \item Compute the angle of refraction at the first face.
  \item The refracted ray hits the second face from inside, at the same
        angle. Compute the exit angle and compare it with $50^\circ$.
  \item What, then, does a thick window do to the direction and to the
        position of a ray crossing it?
\end{enumerate}
\end{exercise}

\begin{exercise}[$\star\star$]\label{exo:g10:refraction:10}
White light strikes a block of flint glass at $55^\circ$. For this
glass, $n_{\text{red}} = 1.61$ and $n_{\text{blue}} = 1.66$. Compute the
refraction angles of the red and blue rays and the angle between them.
Which color ends up closer to the normal?
\end{exercise}

\begin{exercise}[$\star\star$]\label{exo:g10:refraction:11}
An optical fiber has a core of index $1.48$ and a cladding of index
$1.46$.
\begin{enumerate}
  \item Compute the critical angle on the core--cladding boundary.
  \item A guided ray bounces along at exactly the critical angle. Show
        that over \qty{1.0}{km} of straight fiber it travels about
        \qty{14}{m} more than a ray going straight down the axis.
\end{enumerate}
\end{exercise}

\begin{exercise}[$\star\star\star$]\label{exo:g10:refraction:12}
A diver's eye is \qty{2.0}{m} below a perfectly calm surface. Because of
refraction, the diver sees the \emph{entire} sky compressed into a
bright disc directly overhead --- Snell's window.
\begin{enumerate}
  \item Explain why light from the horizon reaches the diver's eye at
        the critical angle of water, $48.8^\circ$ from the vertical.
  \item Compute the radius of the bright disc on the surface, using the
        tangent in the right triangle eye--vertical--disc edge.
  \item What does the diver see on the surface \emph{outside} the disc?
\end{enumerate}
\end{exercise}

\begin{exercise}[$\star\star\star$]\label{exo:g10:refraction:13}
A prism has angles $30^\circ$, $60^\circ$, $90^\circ$. A thin beam of
white light enters perpendicular to one of the sides adjacent to the right angle,
crosses undeviated, and strikes the long face from inside at
$i_1 = 30^\circ$. For this glass, $n_{\text{red}} = 1.51$ and
$n_{\text{blue}} = 1.53$.
\begin{enumerate}
  \item Compute the exit angles of the red and blue rays and the angular
        width of the fan between them.
  \item The prism is replaced by a $45^\circ$--$45^\circ$--$90^\circ$
        prism, so the internal angle of incidence becomes $45^\circ$.
        Show that no light exits through the long face at all.
\end{enumerate}
\end{exercise}

\begin{exercise}[$\star\star\star$]\label{exo:g10:refraction:14}
Periscopes fold light with $45^\circ$--$45^\circ$--$90^\circ$ glass
prisms ($n = 1.50$): the beam enters one short face perpendicularly and
hits the long face from inside at $45^\circ$.
\begin{enumerate}
  \item Show that the long face acts as a perfect mirror. Why is this
        better than a metallic mirror, which absorbs a few percent at
        each reflection?
  \item The periscope floods and the long face is now in contact with
        water. Compute the new critical angle of the glass--water
        boundary and show the mirror fails.
  \item At what angle does the light of question 2 leave the prism into
        the water?
\end{enumerate}
\end{exercise}

\begin{exercise}[$\star\star\star$]\label{exo:g10:refraction:15}
A transatlantic cable runs \qty{6200}{km} of fiber (core index $1.47$)
between London and New York.
\begin{enumerate}
  \item Compute the speed of light in the core and the one-way travel
        time of a pulse.
  \item A radio wave travels at $c$. How long would it need over the
        same distance, and how much round-trip time does the fiber's
        glass cost compared with radio?
  \item Some financial firms pay fortunes to shave milliseconds off this
        link. Propose two distinct physical ways to make the light
        arrive earlier.
\end{enumerate}
\end{exercise}

\section{Problem: The diamond, the pool and the fiber}

\begin{problem}[{Weekend problem --- three lives of Snell's law: a pool
that lies about its depth, a stone cut so that light cannot leave, and
the sixty-millisecond ocean under the internet}]
\label{pb:g10:refraction:1}
One equality of two sines, three careers. In a swimming pool it cheats
your eyes about depth and squeezes the whole sky into a circle. In a
diamond, pushed to its breaking point, it refuses to let light out ---
and jewelers have cut stones around that refusal for centuries. Under
the Atlantic it guides laser pulses inside a strand of glass, and sets
the speed limit of the global internet. This problem works through the
three lives in order, with the same law each time.

\medskip
\noindent\textbf{Part I --- Warming up.}
\begin{enumerate}
  \item Compute the speed of light in water ($n = 1.33$) and in diamond
        ($n = 2.42$). By what factor does diamond slow light down?
  \item A ray in air strikes a pond at $i_1 = 40^\circ$. Compute the
        angle of refraction.
  \item A ray in the water hits the surface from below at $28.9^\circ$.
        At what angle does it exit into the air? State the principle
        this illustrates.
  \item Compute the critical angle of the water--air boundary, and
        describe what happens to an underwater ray hitting the surface
        at $60^\circ$.
  \item The index of air is $1.0003$. Compute the speed of light in air
        and justify the approximation $n_{\text{air}} = 1.00$ used
        throughout.
\end{enumerate}

\medskip
\noindent\textbf{Part II --- The pool.}
\begin{enumerate}[resume]
  \item A pool is \qty{3.0}{m} deep. Using the apparent-depth formula
        (\cref{prop:g10:refraction:depth}), compute the depth a person
        standing above perceives.
  \item Explain with a two-ray sketch (described in words) why a coin
        on the bottom seems to rise: what do rays leaving the water do,
        and where does the eye place their intersection?
  \item A diver's eye is at \qty{3.0}{m}. Compute the radius of Snell's
        window --- the bright disc through which the diver sees the
        whole sky (\cref{exo:g10:refraction:12}).
  \item A lamp on the bottom sends rays upward at $30^\circ$,
        $45^\circ$ and $50^\circ$ from the vertical. For each ray, say
        whether it escapes and at what angle, or what happens instead.
  \item A spear-fisher on the bank aims at a fish. Should the strike be
        aimed above, at, or below the fish's image --- and does the
        fish, looking back, see the fisher displaced too?
\end{enumerate}

\medskip
\noindent\textbf{Part III --- The diamond.}
\begin{enumerate}[resume]
  \item Compute the critical angle of diamond ($n = 2.42$) toward air,
        and compare it with that of glass ($n = 1.52$),
        which is $41.1^\circ$.
  \item A ray inside the stone hits a facet at $30^\circ$ from its
        normal. What happens in the diamond? In the glass imitation?
        Explain in one sentence why the cut of a diamond, combined with
        its tiny critical angle, produces the sparkle.
  \item Diamond is also strongly dispersive: $n_{\text{red}} = 2.41$,
        $n_{\text{blue}} = 2.45$. White light enters a facet at
        $45^\circ$; compute the refraction angle for each color and the
        angular split.
  \item Later, red and blue rays both reach an exit facet from inside
        at $17.0^\circ$. Compute the two exit angles and the split of
        the emerging fan. Compare with the split at entry: what happens
        to dispersion near the critical angle?
  \item A drop of water ($n = 1.33$) coats a facet. Compute the new
        critical angle of the diamond--water boundary, decide the fate
        of the $30^\circ$ ray of question 12, and explain why jewelers
        keep diamonds scrupulously clean.
\end{enumerate}

\medskip
\noindent\textbf{Part IV --- The fiber.}
\begin{enumerate}[resume]
  \item A fiber has a core of index $1.48$ and a cladding of index
        $1.46$. Compute the critical angle on the core--cladding
        boundary.
  \item A bare glass thread in air would have a critical angle of
        $42.5^\circ$ --- seemingly far better. Give two reasons the
        cladding is used anyway. (Think about dust, scratches, and what
        touching the surface would do to total internal reflection.)
  \item Compute the speed of light in the core, then the one-way travel
        time of a pulse along \qty{6000}{km} of fiber.
  \item A ray bouncing at exactly the critical angle travels
        $1/\sin i_c$ times the length of the axis. Compute the extra
        distance over the \qty{6000}{km} and the extra time it costs.
  \item The punchline. A modern fiber pair carries about $10^{13}$ bits
        per second, and this book is about $2.4 \times 10^8$ bits.
        Compute the round-trip (``ping'') time across the ocean and the
        number of copies of this book the fiber moves per second ---
        then state, in one sentence, what Snell's law does for the
        internet.
\end{enumerate}
\end{problem}
